% Physical unit 033; accepted Session 03 after fresh 400-dpi authority replay.
\GordanUnitStart{033}{printed page 28}
\GordanTopFolio{28}
\begingroup\small
\sect{5}

Ad computandum $K$ utimur formula:
\[
 \sqrt{2K\over\pi}=1+2q \qquad \text{fund. p. 184}
\]
vel
\[
 K={\pi\over2}\{1+2q\}^2.
\]
ad computandum $E$ formulis:
\[
\begin{array}{rcll}
 {2K\over\pi}{2K\over\pi}-{2K\over\pi}\cdot {2E\over\pi}&=&8q+16q^2+32q^3\ldots, &135\vspace{0.2em}\\
 {2K\over\pi}\cdot{2K\over\pi}&=&1+8q+24q^2+32q^3\ldots, &105
\end{array}
\]
Subtrahendo fit:
\[
 {2K\over\pi}\cdot {2E\over\pi}=1+8q^2
\]
et dividendo per aequationem superiorem:
\[
 {2E\over\pi}={1+8q^2\over (1+2q)^2}.
\]
Haec est altera computandi $\Sigma(l_0)$ ratio.

Est enim:
\[
 \Sigma(l_0)={b\over x'}E \qquad \text{Cap. VI \S\ 3}
\]
unde fit:
\[
 {2\over\pi}\Sigma(l_0)=b{1+8q^2\over \{1-2q\}^2}.
\]

\sect{6}

Restat, ut $v$ computetur; quod quoniam explicitum non existit, relatas ad id functiones trigonometricas fractionis
\[
 {iv\pi\over2K}
\]
computare placet.

Proinde novam constantem $A$ inducere volumus, definiendam aequatione
\[
 \sin A\cos {iv\pi\over2K}=1
\]
vel:
\[
 \sin {iv\pi\over2K}=i\cot A
 \qquad
 \tg {iv\pi\over2K}=i\cos A
\]
e quibus intelligitur, $A$ tantummodo computari debere.
\endgroup
